%% Document created 28 April 2021 automatically 
%% from /Users/massimosotgia/Desktop/uni_at_DIFI/Lab_C03/setup.py 

%% Copyright (C) Mattia Sotgia et al. 2021
%% Using class lab_unige.cls
%                                                            
%                                                            
%   **                 **             ******   ****   ****   
%  /**                /**            **////** *///** */// *  
%  /**        ******  /**           **    // /*  */*/    /*  
%  /**       ´´´´´´** /******      /**       /* * /*   ***   
%  /**        ******* /**///**     /**       /**  /*  /// *  
%  /**       **´´´´** /**  /** **  //**    **/*   /* *   /*  
%  /********//********/****** /**   //****** / **** / ****   
%  ////////  //////// /////   //     //////   ////   /´///   
%                                                            
%                                                            
\documentclass[italian, a4paper, 10pt, twocolumn]{../../style/lab_unige}
\usepackage[a4paper, margin=1.75cm, footskip=0.25in]{geometry}

\usepackage[utf8]{inputenc}
\usepackage[T1]{fontenc}

\usepackage[italian]{babel}

% \usepackage{biblatex}

\usepackage[bookmarksopen=true, citebordercolor={0 1 0}, linkbordercolor={1 0 0}, urlbordercolor={0 1 1}]{hyperref}
\usepackage[numbered]{bookmark}

\usepackage{graphicx}
\graphicspath{{../fig/}}
\usepackage{array}
\usepackage{tabulary}
\usepackage{booktabs}

% FOUNDAMENTAL
\usepackage{../../style/custom}

\usepackage{physics}

\usepackage{breqn}
\usepackage{cuted}
\usepackage{txfonts}

\usepackage{lipsum}

%% Define ref types
\newcommand{\reftab}[1]{Tabella {\ref{#1}}}%
\newcommand{\reffig}[1]{Figura {\ref{#1}}}%
\newcommand{\refeqn}[1]{Eq. ({\ref{#1}})}%
\newcommand{\ChiSqr}{$\chi^2$\space}
\newcommand{\ChiNdf}{$\chi^2/\text{ndf}$}
\newcommand{\cernroot}{\texttt{root}}
\newcommand{\treSigma}{$3\sigma$}
\newcommand{\stdErr}[1]{$\varepsilon_{#1}$}
\newcommand{\mstdErr}[1]{\varepsilon_{#1}}
\newcommand{\gLab}{$g_t~=~(~9.8056~\pm~0.0001~\text{ stat}~)~\text{ m/s}^2$}
%% PAPER ONLY custom Macros


%%
\setlength{\columnsep}{6mm}

\begin{document}
    \twocolumn[
    \begin{@twocolumnfalse}
        \title{
            {\raggedright \includegraphics[width=0.2\linewidth]{../../style/lab_mark.pdf}\\}
            Metodo Volt-Amperometrico E Circuito RC 
        }
        \author{
        Eugenio Dormicchi\textsuperscript{1, 2},
        % Riccardo Pizzimbone\textsuperscript{1}, 
        Giovanni Oliveri\textsuperscript{1},
        Mattia Sotgia\textsuperscript{1}
        }

        \date{
            \textit{
            \textsuperscript{1}Gruppo C03, Esperienza di laboratorio n. 8 \\
            \textsuperscript{2}In presenza in laboratorio per la presa dati\\}
            % Università degli Studi di Genova, Dipartimento di Fisica.\\
            \vspace{2ex}
            (Presa dati
            28 Aprile 2021, 15:00– 18:00; Analisi dati 
            4 Maggio 2021)
            
        }
        \maketitle
        
        \begin{abstract}
            \textit{Obiettivo-- }
            
            \textit{Metodi-- }
            Circuiti elettrici, carica e scarica condensatori, Fit polinomiali ed esponenziali. 
            \textit{Risultati-- }
        
            \textit{Conclusione-- }
        
        
        \end{abstract}
        \vspace{2em}
    \end{@twocolumnfalse}
    ]

    %%%% CORPO DEL TESTO
    %%%% CORPO DEL TESTO

    \section{Introduzione}
    \label{section:introduction}

    % \section{Strumentazione}
    % \label{section:strument}

    \section{Metodi}
    \label{section:methods}

    % \section{Analisi dati}
    % \label{section:analysis}

    \section{Risultati}
    \label{section:results}

    \section{Conclusione}
    \label{section:conclusion}

    \subsection{Controlli}

    \subsection{Possibili errori sistematici}

    \appendix

    \setcounter{table}{0}
    \renewcommand{\thetable}{A\arabic{table}}

    \section{Dati estesi}
    

\end{document}
    
